\section*{Capítulo \ref{ch:b3:cell-cycle-apoptosis} --- Controle do ciclo celular e morte celular programada}

\begin{solution}{exo:b3:cell-cycle-apoptosis:1}
G1, S, G2, M. Progressão em G1 e \omterm{def:b3:cell-cycle-apoptosis:checkpoints}{ponto de restrição}: \omterm{def:b3:cell-cycle-apoptosis:cdk}{ciclina}
D--Cdk4/6; G1/S: \omterm{def:b3:cell-cycle-apoptosis:cdk}{ciclina} E--Cdk2; S: \omterm{def:b3:cell-cycle-apoptosis:cdk}{ciclina} A--Cdk2; G2/M: \omterm{def:b3:cell-cycle-apoptosis:cdk}{ciclina}
B--Cdk1. O \omterm{def:b3:cell-cycle-apoptosis:cdk}{complexo promotor da anáfase} (APC/C) destrói a \omterm{def:b3:cell-cycle-apoptosis:cdk}{ciclina} B e a
securina e encerra a mitose.
\end{solution}

\begin{solution}{exo:b3:cell-cycle-apoptosis:2}
Hartwell: os mutantes \emph{cdc} da levedura de brotamento, parados em
etapas definidas, e o conceito de Start. Nurse: o \emph{cdc2} da
levedura de fissão como a cinase que marca o tempo da mitose, e seu
homólogo humano (Cdk1) resgatando a levedura --- a conservação. Hunt: a
\omterm{def:b3:cell-cycle-apoptosis:cdk}{ciclina} em óvulos de ouriço-do-mar, uma proteína sintetizada ao longo
do ciclo e destruída a cada divisão.
\end{solution}

\begin{solution}{exo:b3:cell-cycle-apoptosis:3}
\omterm{def:b3:cell-cycle-apoptosis:checkpoints}{Ponto de restrição} (fim de G1): fatores de crescimento, nutrientes,
tamanho; age sobre a \omterm{def:b3:cell-cycle-apoptosis:cdk}{ciclina} D--Cdk4/6 e sobre a Rb. G2/M: DNA intacto
e replicado; a \omterm{def:b3:dna-repair:ddr}{ATM}/ATR inibem a Cdc25, mantendo a Cdk1 fosforilada e
inativa. Montagem do fuso: cada cinetócoro ligado; cinetócoros não
ligados produzem complexos Mad2--Cdc20 que inibem o APC/C.
\end{solution}

\begin{solution}{exo:b3:cell-cycle-apoptosis:4}
\omterm{def:b3:cell-cycle-apoptosis:apoptosis}{Apoptose}: a célula encolhe, a \omterm{def:b3:chromatin-epigenetics:nucleosome}{cromatina} se condensa, o DNA é cortado
numa escada, a membrana forma bolhas, mas permanece selada, a
\omterm{def:b3:cell-cycle-apoptosis:apoptosis}{fosfatidilserina} fica exposta, o corpo é englobado, não há inflamação.
\omterm{def:b3:cell-cycle-apoptosis:apoptosis}{Necrose}: a célula incha, a membrana se rompe, o conteúdo vaza, o DNA é
degradado ao acaso, há inflamação. Executoras: as \omterm{def:b3:cell-cycle-apoptosis:apoptosis}{caspases}, proteases
de cisteína que cortam depois de um aspartato.
\end{solution}

\begin{solution}{exo:b3:cell-cycle-apoptosis:5}
Subida de $5$ a $30$ a $1$ por minuto: \qty{25}{min}. Degradação de
$30$ a $5$ com meia-vida de \qty{2}{min}: $\log_{2}(30/5)\times 2 =
\qty{5.2}{min}$. Período de cerca de \qty{30}{min}, com a mitose
ocupando \qty{17}{\%} dele.
\end{solution}

\begin{solution}{exo:b3:cell-cycle-apoptosis:6}
(a) $C$ nunca pode cair abaixo de $C_{1}$: travada em mitose, no ramo
alto. (b) Sem a fosforilação inibitória: o limiar $C_{2}$ baixa, e a
mitose começa cedo, num tamanho pequeno (o fenótipo ``wee''). (c) A
Cdk1 permanece fosforilada: o ramo alto fica inalcançável, com parada
em G2 e células alongadas. (d) Como em (b), e o \omterm{def:b3:cell-cycle-apoptosis:checkpoints}{ponto de checagem} G2/M,
que age por meio de Wee1/Cdc25, já não consegue segurar a célula.
\end{solution}

\begin{solution}{exo:b3:cell-cycle-apoptosis:7}
Núcleo em G1 em citoplasma de fase S: suas origens estão licenciadas e
o citoplasma fornece \omterm{def:b3:cell-cycle-apoptosis:cdk}{ciclina} E/A--Cdk2 ativa, de modo que ele entra em
S de imediato. Núcleo em G2 em citoplasma de fase S: suas origens já
dispararam e não foram relicenciadas (os fatores de licenciamento foram
destruídos ou inibidos pelas CDKs), de modo que ele não pode replicar
de novo; espera até que o parceiro chegue a G2 e os dois entrem juntos
em mitose.
\end{solution}

\begin{solution}{exo:b3:cell-cycle-apoptosis:8}
O cinetócoro não ligado é um catalisador: ele converte continuamente a
Mad2 em sua conformação ativa, produzindo um inibidor difusível da
Cdc20 que se espalha pela célula e mantém o APC/C desligado em toda
parte --- um único cinetócoro basta. Sem a Mad2, o APC/C é ativado
assim que a Cdk1 o liga, estejam ou não os cromossomos ligados: a
anáfase começa com cromossomos soltos, as filhas ficam aneuploides, e o
organismo (ou a linhagem celular) morre de perda de cromossomos.
\end{solution}

\begin{solution}{exo:b3:cell-cycle-apoptosis:9}
Sem \emph{ced-9}: células que deveriam viver morrem --- letalidade
embrionária por morte demais. Sem \emph{ced-3}: nenhuma das $131$
mortes ocorre, as células sobrevivem e se diferenciam, e o verme é
viável. Sem os dois: nenhuma morte --- o fenótipo de \emph{ced-3}. Como
remover a executora cancela o efeito de remover a guardiã, a CED-9 age
a montante, mantendo CED-4/CED-3 inativas; quando ela se vai, elas
ficam ativas, a menos que também estejam ausentes.
\end{solution}

\begin{solution}{exo:b3:cell-cycle-apoptosis:10}
Com $A = aC$, $\mathrm{d}C/\mathrm{d}t = k_{s} - k_{d}aC^{2}$, uma
equação de uma variável com estado estacionário
$C^{*} = \sqrt{k_{s}/(k_{d}a)}$ e inclinação $-2k_{d}aC^{*} < 0$ ali:
qualquer desvio decai monotonicamente, e uma única variável de primeira
ordem não pode oscilar. A curva em S dá dois ramos estáveis separados
por uma zona proibida, de modo que o estado precisa saltar e não pode
se assentar; um atraso explícito na retroalimentação daria oscilações
por outra rota --- o freio chegando tarde demais, ultrapassando o ponto
e assim por diante ---, como em muitos ritmos hormonais.
\end{solution}

\begin{solution}{exo:b3:cell-cycle-apoptosis:11}
Cada sensora que chega é capturada por uma guardiã, de modo que nada
acontece até que todas as $N$ guardiãs estejam ocupadas, em
$t \approx N/r$; as sensoras seguintes ficam livres, ativam Bax/Bak,
cuja formação de poros sobe abruptamente com o número delas, e as
mitocôndrias se abrem em minutos. Um estímulo mais forte encurta o
tempo apenas por meio de $r$; acima do limiar o desfecho é o mesmo. O
venetoclax ocupa guardiãs, reduzindo o $N$ efetivo: o tempo encolhe, e
uma célula preparada --- já perto de $N$ --- morre de imediato.
\end{solution}

\begin{solution}{exo:b3:cell-cycle-apoptosis:12}
O excesso de Bcl-2 bloqueia a \omterm{def:b3:cell-cycle-apoptosis:pathways}{via intrínseca}: os linfócitos~B que
deveriam morrer quando seus sinais de crescimento acabam sobrevivem, e
a população cresce por falha de morte, à taxa lenta com que tais
células são feitas --- indolente, até que uma segunda lesão acrescente
proliferação. A perda da Rb remove o \omterm{def:b3:cell-cycle-apoptosis:checkpoints}{ponto de restrição}: os precursores
da retina percorrem o ciclo sem fatores de crescimento e se dividem tão
depressa quanto o motor permite --- agressivo. Os dois tumores são os
dois programas falhando cada um por si: um controle de morte e um
controle de divisão, bastando a perda de qualquer um deles para um
tumor, o segundo mais rápido.
\end{solution}

\begin{solution}{pb:b3:cell-cycle-apoptosis:1}
\textbf{1.} $10^{7}\times 3500/5 = 7\times 10^{9}$ células por dia,
cerca de \num{80000} por segundo.
\textbf{2.} $3500/5 = 700$ por cripta por dia; $22\times 2^{5} = 704$.
\textbf{3.} $80\times 365 \approx \num{29000}$ divisões; $\times 0.6
\approx \num{17500}$ mutações.
\textbf{4.} As células de trânsito são descartadas em dias, levando
consigo suas mutações; a célula-tronco persiste a vida toda, de modo
que cada uma de suas mutações é conservada e cada divisão acrescenta
mais --- ciclar devagar as minimiza.
\textbf{5.} A vilosidade se esvazia nos cerca de $5$ dias de seu
trânsito normal; reconstruí-la leva um dia de recuperação, cinco
divisões a cada \qty{12}{h} e a migração vilosidade acima --- quatro a
cinco dias, durante os quais a barreira fica nua: as divisões rápidas
de trânsito do intestino fazem dele uma das primeiras vítimas da
radiação.
\textbf{6.} Um corpo apoptótico selado não libera conteúdo algum num
lúmen cheio de bactérias e não provoca inflamação; a \omterm{def:b3:cell-cycle-apoptosis:apoptosis}{necrose} derramaria
enzimas e sinais e inflamaria a mucosa a cada célula descartada.
\textbf{7.} $(40 - 8)/2 = \qty{16}{h}$.
\textbf{8.} $\log_{2}(40/8)\times \qty{10}{min} = 2.3\times 10 =
\qty{23}{min}$.
\textbf{9.} Período de cerca de \qty{16.4}{h}; a Cdk1 fica ativa
\qty{2.3}{\%} do tempo.
\textbf{10.} Não é o $k_{s}$ sozinho, mas o tempo em que a célula é
segurada nos limiares: a célula somática espera no \omterm{def:b3:cell-cycle-apoptosis:checkpoints}{ponto de restrição}
por fatores de crescimento e em G2/M por seu DNA, e a célula-tronco
espera mais que a de trânsito. O óvulo de rã não tem \omterm{def:b3:cell-cycle-apoptosis:checkpoints}{pontos de checagem}
e tem um citoplasma abastecido de tudo, de modo que só o oscilador nu
fixa seu período.
\textbf{11.} No ramo baixo, num nível de \omterm{def:b3:cell-cycle-apoptosis:cdk}{ciclina} acima do $C_{2}$
normal: o \omterm{def:b3:cell-cycle-apoptosis:checkpoints}{ponto de checagem} deslocou a curva em S para a direita ao
inibir a Cdc25, e por isso o salto não ocorre. Com o reparo, a Cdc25 é
liberada, o limiar cai abaixo da \omterm{def:b3:cell-cycle-apoptosis:cdk}{ciclina} acumulada, e a célula entra em
mitose de imediato --- razão pela qual células liberadas de um \omterm{def:b3:cell-cycle-apoptosis:checkpoints}{ponto de
checagem} entram em mitose de modo sincronizado.
\textbf{12.} Nem a \omterm{def:b3:cell-cycle-apoptosis:cdk}{ciclina} B nem a securina são degradadas: a célula
permanece em metáfase, com a coesina intacta e a Cdk1 ativa,
indefinidamente; ela morre ali ou acaba escapando com um genoma não
segregado.
\textbf{13.} $10^{11}$ divisões por dia.
\textbf{14.} Com o \omterm{def:b3:cell-cycle-apoptosis:checkpoints}{ponto de checagem}, $10^{11}/5\times 10^{4} = 2\times 10^{6}$ filhas aneuploides por dia; sem ele, $2\times 10^{9}$.
\textbf{15.} $2\times 10^{4}$ por dia; ao longo de $80$ anos,
$6\times 10^{8}$.
\textbf{16.} $10^{9}/50 = 2\times 10^{7}$ más segregações por dia: todo
dia o tumor tenta vinte milhões de cariótipos novos, e a seleção guarda
os que crescem mais depressa ou resistem a um fármaco.
\textbf{17.} Com cada divisão segregando mal, nenhuma linhagem celular
embrionária conserva um genoma euploide e o embrião morre. A perda
parcial dá aneuploidia ocasional em células que a sobrevivem (sobretudo
sem p53), uma instabilidade cromossômica que fornece variação a um
tumor sem matar o organismo.
\textbf{18.} Um erro meiótico põe o cromossomo extra no gameta, de modo
que o zigoto e cada célula dele derivada são trissômicos; um erro
mitótico ocorre numa célula somática e é herdado apenas por seu clone.
\textbf{19.} $10^{11}\times \qty{1}{ng} = \qty{100}{g}$ por dia,
\qty{36}{kg} por ano --- cerca de metade da massa do corpo a cada ano.
\textbf{20.} \qty{20}{g} de proteína por dia, cerca de um terço da
ingestão alimentar e alguns por cento da renovação proteica total do
corpo.
\textbf{21.} $10^{11}$ corpos a $24$ por macrófago por dia: $4\times
10^{9}$ macrófagos em tempo integral, \qty{4}{\%} do tempo dos
$10^{11}$ macrófagos.
\textbf{22.} Um corpo lisado libera proteases, DNA, ATP e outros sinais
de alarme que causam inflamação e podem provocar autoimunidade contra
antígenos nucleares; o corpo intacto se anuncia com \omterm{def:b3:cell-cycle-apoptosis:apoptosis}{fosfatidilserina} em
sua monocamada externa (e pela perda de seus sinais de ``não me coma'')
e atrai fagócitos com nucleotídeos liberados.
\textbf{23.} A \omterm{def:b3:cell-cycle-apoptosis:pathways}{via intrínseca} está bloqueada na mitocôndria. A rota
extrínseca ainda funciona onde a caspase-8 ativa a caspase-3
diretamente, embora sua amplificação pela Bid se perca. As células se
acumulam por não morrer, e não por se dividirem mais depressa ---
crescimento lento ---, até que uma lesão posterior (muitas vezes o
\emph{MYC}) acrescente proliferação.
\textbf{24.} A \omterm{def:b3:cell-cycle-apoptosis:apoptosis}{apoptose} leva horas, precisa de ATP e passa por etapas
que podem ser bloqueadas --- inibidores de \omterm{def:b3:cell-cycle-apoptosis:apoptosis}{caspase}, o limiar
mitocondrial ---, de modo que os neurônios da penumbra podem ser
resgatados se tratados a tempo; os neurônios do núcleo morrem de falha
energética em minutos, sem programa algum a interromper.
\textbf{25.} Período do oscilador nu de cerca de \qty{16}{h}; uns
$2\times
10^{4}$ células aneuploides em divisão por dia num corpo
normal; \qty{100}{g} de células recicladas por \omterm{def:b3:cell-cycle-apoptosis:apoptosis}{apoptose} a cada dia.
\end{solution}
